When the sources retrieved for a query disagree, the right response depends on \emph{why} they disagree: a
question with several valid answers, a contested topic, stale information and a false source each call for
different behaviour. \citet{cattan2025dragged} showed that telling a generator the conflict type improves its
responses by about 24 points, but that predicting the type is hard even for frontier LLMs. We study whether a
\emph{typed decision model}, which returns calibrated probabilities over fixed options in a single forward pass
instead of generating text, can do this classification fast and cheaply. Using the \conflicts benchmark and the
original paper's prompt verbatim, we evaluate the decision model \jev against Gemini models on identical inputs.
Asked directly, \jev reaches \jevacc\% accuracy, statistically tied with Gemini~2.5 Flash, the paper's model
(\gtwofiveacc\%, $p\jevvstwofivep$), at \costratiotwofive$\times$ lower cost and \latratiotwofive$\times$ lower
latency, and it is the fastest of the \nconfigs model configurations we test. It trails the current Gemini~3.8 Flash by
\jevvsthreeeightdiffabs points; a confidence-gated cascade that sends \jev's least certain quarter of queries to that
LLM recovers most of the gap (\cascflattfacc\% vs.\ \gthreeeightacc\%, $p\cascflattfp$) at \cascflattfcost\% of
its cost. Decomposing the taxonomy into yes/no questions fails (\treeoneacc\%) because of a yes-bias, and a repair
chosen on development data does not transfer to other decision models. Pre-specified diagnostics show that
\jev's documented weakness with dates does not matter for this task, and that its sensitivity to source order and
its ability to notice a planted false source match the LLM's. We release code, prompts and all model outputs.
